\section{Related Work}\label{sec:related}

The Introduction set out the thesis question: whether a general-purpose zkVM can verify an algebraic-hash post-quantum signature feasibly on consumer hardware, and where the deployability wall for post-quantum anonymity then sits. Five strands of prior work bear on it. Post-quantum signatures and anonymous credentials establish what BDEC is assembled from; transparent succinct arguments and zero-knowledge virtual machines are the two substrates this thesis sets against each other; and precompile-based cost reduction is the instrument that makes the zkVM substrate measurable at all. We read each strand with one question in view: what the existing evaluation literature measures and what it leaves unmeasured. The closing subsection turns to the omission this thesis fills. A strand is included when it changes the comparison on one of four axes, namely post-quantum credential security, in-circuit post-quantum signature verification, zkVM proving cost, or update churn, which is why deployed verifiable-credential standards and BBS appear as background rather than as direct baselines.

\subsection{Post-Quantum Signatures}

The threat that large-scale quantum computers pose to RSA- and discrete-logarithm-based signatures has driven a multi-year standardisation effort. The NIST post-quantum cryptography process has produced two general-purpose signature standards, ML-DSA~\cite{fips204} (a lattice-based scheme derived from Dilithium) and SLH-DSA~\cite{fips205} (a stateless hash-based scheme derived from SPHINCS+), with a third, Falcon (FN-DSA, draft FIPS~206~\cite{fips206}), also lattice-based.
These signatures are designed for general-purpose authentication and are well suited to consumer hardware; they are, however, not well suited as the underlying primitive of an anonymous-credential system because their verification predicates do not embed cheaply inside a zkSNARK. Lattice-based verification involves expensive rejection sampling and norm checks, and hash-based verification involves on the order of $10^5$ symmetric-primitive invocations per signature~\cite{fips205}, both of which translate into large in-SNARK constraint counts. This concerns the \emph{general} verification predicate. Narrower-functionality results, such as post-quantum Privacy Pass~\cite{policharla2023pqprivacypass} and hardware or precompile acceleration of the hash-based families, reduce these costs in restricted settings, so the ``not well suited'' judgement applies to the general anonymous-credential predicate and admits these exceptions in restricted ones.

A complementary line of work develops post-quantum signatures whose verification predicate is deliberately structured for a low constraint count when encoded as an R1CS. Picnic, Banquet, and Limbo build signatures around the MPC-in-the-head paradigm~\cite{chase2017picnic,baum2021banquet,desaintguilhem2021limbo}. Loquat~\cite{cryptoeprint:2024/868} departs from this line: it bases unforgeability on the Legendre PRF, a symmetric-primitive assumption distinct from lattice and code-based hardness, and is designed so that its verification circuit fits into a comparatively small R1CS. PLUM~\cite{10.1007/978-981-95-2961-2_6}, the focus of this thesis, replaces Loquat's Legendre PRF with the $t$-th power residue PRF and replaces the FRI low-degree test with the STIR protocol. The result is, at 128-bit security, signatures approximately $1.5\times$ smaller, verification approximately $4\times$ faster, and approximately $1.28\times$ fewer R1CS constraints per verification, while preserving Loquat's security guarantees~\cite{10.1007/978-981-95-2961-2_6}. This trio of improvements is what motivates evaluating PLUM as a replacement for Loquat in the BDEC framework; whether the replacement is strictly favourable inside a zkVM, as opposed to the static-circuit setting these figures describe, is one of the questions this thesis answers. A direct competitor in this SNARK-friendly line is CAPSS~\cite{feneuil2025capss}, a framework for SNARK-friendly post-quantum signatures; the schemes differ as much in \emph{assumption} as in constraint count, since Loquat and PLUM rest on the Legendre and power-residue PRFs, MPC-in-the-head designs on generic symmetric primitives, lattice schemes on SIS/LWE, and every algebraic-hash argument additionally on a random-oracle instantiation of its hash, so a constraint-count comparison alone does not settle the choice.

\subsection{Anonymous Credentials}

Anonymous credentials originate with Chaum's pseudonymous-credentials proposal~\cite{chaum1985}, formalised by Camenisch and Lysyanskaya (CL)~\cite{camenisch2004signatures}, and have since developed into a substantial line including Pointcheval--Sanders signatures, Coconut, and the BBS family~\cite{pointchevalsanders2016,sonnino2019coconut,bbs04}. These remain the basis of deployed selective-disclosure systems (W3C Verifiable Credentials~\cite{w3cvc}, BBS-based showing), which we treat as deployment baselines rather than post-quantum-anonymous-credential baselines, since none is quantum-resistant.
These classical constructions rely on pairing-based or discrete-logarithm-based assumptions, and are therefore not quantum-resistant. The post-quantum case has, until recently, been pursued largely through lattice-based primitives: Jeudy et al.\ at Crypto~2023~\cite{jeudy2023lattice} provide an efficient lattice-based anonymous-credential scheme, leveraging the SIS and LWE assumptions; subsequent work has refined and applied this template. The known limitations of the lattice approach include large proof and key sizes and, in the Jeudy et al.\ scheme as published, support for only a single message per credential.

The Blockchain-based Digital Education Credential (BDEC) framework of Li, Zhang et al.~\cite{10.1007/978-981-96-0957-4_3} takes a different path. BDEC is generic over a signature scheme and a zkSNARK, and the authors instantiate it with Loquat and Aurora to obtain a post-quantum anonymous-credential system whose security reduces to collision-resistance and the Legendre PRF. The system supports unforgeability, anonymity, unlinkability, conditional linkability, and revocability, and the authors report concrete-performance improvements over the lattice baseline of Jeudy et al.~\cite{jeudy2023lattice}. The BDEC construction is the anonymous-credential framework this thesis builds on; we examine it in detail in Section~\ref{sec:prelim}. We distinguish BDEC-as-proposed from the instantiation we measure: BDEC's anonymity and unlinkability require the underlying argument to be zero-knowledge, which the consumer-hardware zkVM receipt we produce does not yet satisfy (Section~\ref{sec:security}), so our artifact does not inherit those properties unconditionally.

\subsection{Transparent Succinct Arguments and Low-Degree Testing}

A zkSNARK proof system consists of an interactive-oracle-proof (IOP) compiled into a non-interactive argument via a hash-based transformation, most commonly the Ben-Sasson--Chiesa--Spooner (BCS) transformation~\cite{bensasson2016iop}.
The BCS transformation requires a low-degree test that confirms a committed function is close to a low-degree polynomial. Two such tests dominate recent transparent-argument constructions: FRI (Fast Reed--Solomon Interactive Oracle Proof of Proximity)~\cite{bensasson2018fri} and STIR (Reed--Solomon proximity testing with fewer queries)~\cite{cryptoeprint:2024/390}. STIR achieves lower query complexity and a smaller code-rate after folding, at the cost of slightly more involved verifier logic; both rest on Reed--Solomon proximity-gap results~\cite{cryptoeprint:2020/654}.

Aurora~\cite{cryptoeprint:2018/828} is the canonical transparent succinct argument for R1CS in the IOP-plus-BCS family. It produces zero-knowledge proofs of R1CS satisfiability with polylogarithmic proof size but \emph{linear} verifier complexity, where the verifier reads the explicit R1CS matrices~\cite[Thm.~1.2]{cryptoeprint:2018/828}, with no relation-specific trusted setup.
 (The holographic variant Fractal~\cite{cryptoeprint:2019/1076} attains a polylogarithmic verifier through preprocessing; we use plain Aurora as the BDEC baseline.) Aurora is the zkSNARK choice of the BDEC framework as instantiated by Li, Zhang et al.~\cite{10.1007/978-981-96-0957-4_3}, and is therefore the static-circuit baseline against which we compare in this thesis. Other transparent succinct arguments share the property of relation-specific R1CS encoding and therefore the same cryptographic rigidity discussed in Section~\ref{sec:intro}; the holographic ones (Fractal, Marlin~\cite{chiesa2020marlin}) differ in one respect that matters for the recompile cost we measure, namely that they preprocess each relation into an index, so their per-circuit setup is large rather than negligible, unlike the non-preprocessing Aurora that the BDEC baseline uses.

\subsection{Zero-Knowledge Virtual Machines}

A zero-knowledge virtual machine (zkVM) is a zkSNARK whose underlying relation is the correctness of executing a guest program on a fixed instruction-set architecture~\cite{cryptoeprint:2023/1032,lavin2024surveyapplicationszeroknowledgeproofs}. The term ``zkVM'' is ecosystem terminology: the \emph{default} receipt these systems produce is succinct but not necessarily zero-knowledge, a gap that becomes this thesis's central security finding (Section~\ref{sec:security}). Several zkVM designs are mature enough to be considered for credential-system deployment:

\begin{itemize}
  \item \textbf{RISC Zero}~\cite{bruestle2023risczero} is a RISC-V zkVM whose prover is a recursive STARK over the BabyBear field. Application-defined precompiles are supported via the Zirgen big-integer dialect.
  \item \textbf{SP1}~\cite{succinct2024sp1} is a RISC-V zkVM whose prover is built on the Plonky3 toolkit. SP1 supports application-defined precompiles via its custom-precompile framework, and offers an explicit chained API for wrapping a succinct STARK receipt in Groth16 or PLONK to obtain a zero-knowledge, succinctly-verifiable proof.
  \item \textbf{Jolt}~\cite{cryptoeprint:2023/1217} is a RISC-V zkVM whose prover uses a sum-check-and-lookup approach instead of segment-based STARK proving. Jolt's design optimises for guest cycle counts of moderate size and is included in this discussion for completeness.
\end{itemize}

A separate strand of work, exemplified by zkLLVM~\cite{Kaskov2024zkLLVM}, takes a compiler approach in which high-level programs are compiled into circuit-friendly intermediate representations rather than RISC-V; this approach inherits the cryptographic-rigidity property of static circuits and is therefore positioned closer to the Aurora regime than to the zkVM regime as we use the term in this thesis. The same applies to circuit-frontends such as Circom~\cite{10002421} and Noir, which target a single underlying SNARK system and lock the verification predicate at compile time.

\subsection{Precompile-Based Cost Reduction}

The general technique of replacing emulated computation in a zkVM with a native AIR-based sub-circuit, recognised by the prover via a syscall, has been applied to several primitives. The R0VM~2.0 release adds precompiles for BN254 and BLS12-381 elliptic-curve operations and reports reducing Ethereum block-proving from approximately $35$ minutes to approximately $44$ seconds~\cite{risczero2025r0vm2}. Each such precompile introduces a primitive-specific arithmetisation that becomes part of the trusted computing base and must be audited, as our Griffin precompile is (premise~(T2), Section~\ref{sec:security}). The speedup therefore comes with an audit obligation.
A community RSA precompile for zkVMs has likewise been reported to cut RSA signature verification from tens of millions of guest cycles to the order of $10^5$, roughly two orders of magnitude.\footnote{Reported in vendor/community engineering write-ups; we cite it only as an order-of-magnitude indication of the precompile pattern, not as a peer-reviewed measurement.}
SP1's standard library ships SHA-256 and ECDSA precompiles by default~\cite{succinct2024sp1}; RISC Zero ships SHA-256 and BigInt-256 modular-multiplication precompiles by default~\cite{bruestle2023risczero}.

These results share a structural pattern. A fixed primitive whose software emulation over the prover's native field would dominate the trace is replaced by a native AIR, and the per-call cost from the guest's perspective collapses to a single syscall. The Griffin-$\mathbb{F}_p^{192}$ precompile proposed in this thesis sits within this pattern, with the additional consideration that the primitive in question is parameterised over a non-trivial $192$-bit prime field whose support requires a corresponding BigInt-192 precompile suite for soundness.

\subsection{Evaluation Methodology and the Gap This Work Addresses}

The recent literature on PQ-anonymous-credential systems reports evaluation along the conventional dimensions: prover wall-clock time, verifier wall-clock time, proof size, and, less commonly, peak memory~\cite{10.1007/978-981-96-0957-4_3,10.1007/978-981-95-2961-2_6}. The same is true of the zkVM literature, which adds cycle counts and segment counts to the standard list.

The cost incurred when the verification predicate must be changed, in response to a change in the presentation structure, a primitive substitution, or a security-parameter retuning, is not reported alongside these dimensions. The concept that such changes carry a migration cost is itself not new: crypto-agility work motivates this cost as a design concern, and post-quantum verifiable-credential proposals have noted the need to re-derive parameters when the underlying primitive changes.
For a static-circuit deployment that change requires recompilation of the R1CS matrices and re-derivation of any parameters that depend on the constraint count; for a zkVM deployment it requires only an edit to the guest program. Without accounting for this cost, comparisons between the two regimes omit the regeneration-and-re-audit footprint that each relation change imposes on the static side. We measure that footprint in Section~\ref{sec:evaluation}, and we report there that its wall-clock component is sub-second against a transparent non-preprocessing argument, so the separation between the two regimes is one of re-audit footprint rather than recompile time, and is itself conditional on the static baseline's preprocessing model. This is a deployment-lifecycle cost, paid once per relation change, and it falls outside the per-proof wall-clock figures that the prior work reports. The prior work that names the concern stops short of measuring it. We introduce \emph{update churn} as a measurement dimension and report it alongside the conventional dimensions for both regimes.

The closest prior work confirms that this measurement dimension is unreported. Post-quantum anonymous credentials have been built lattice-natively (Bootle et al.~\cite{BootleLNS23}; Argo et al.~\cite{ArgoGJLRS24}, the latter with sub-$80$~KB credentials and sub-second native showing), classically from zkSNARKs composed with existing identity infrastructure (zk-creds~\cite{rosenberg2023zkcreds}, which is pairing-based and not post-quantum), and on a hand-rolled relation-specific STARK for a lattice signature (Policharla et al.~\cite{policharla2023pqprivacypass}); none runs its verification predicate inside a general-purpose zkVM, and none reports the per-change deployment cost. SNARK-friendly post-quantum signature \emph{design} has likewise been generalised into a framework (CAPSS~\cite{feneuil2025capss}), but as a signature-design methodology, not a measured anonymous-credential-in-zkVM study. On the substrate side, individual post-quantum \emph{signatures} have recently been verified inside general-purpose zkVMs as engineering demonstrations (Falcon and SPHINCS+ in Cairo~\cite{s2morrow2025}; ML-DSA in SP1, Groth16-wrapped~\cite{kota2025dilithiumzk}); we therefore do not claim the first post-quantum signature in a zkVM. Our claim is narrower and measured: to our knowledge, the first reported execution-and-proving study of an algebraic-hash post-quantum \emph{anonymous credential} (BDEC instantiated with PLUM) inside a general-purpose zkVM, with the update-churn dimension and the dual zero-knowledge/post-quantum obstruction that none of the above confronts. We position against the recent zkVM systematisation~\cite{yang2026sokzkvm}, which benchmarks general-purpose zkVMs on classical workloads but measures no post-quantum anonymous credential. None of these works reports the per-change update-churn cost; that is the gap this work fills. The next section develops the primitives the measurement rests on, beginning with PLUM, BDEC, and the two argument substrates.
